\documentclass[12pt]{article}
\usepackage[utf8]{inputenc}
\usepackage[T1]{fontenc}
\usepackage{hyperref}
\title{Minang Language Resources for K-12 Curriculum Development}
\author{Compiled by Claude Sonnet 5.5 (claude-sonnet-5-5) for Andr\'as Kornai}
\date{}
\begin{document}
\maketitle

\section{Introduction}
Minang (Mineng, Menang; ISO 639-3 xrg) is a Noongar variety of the Albany area on the south coast of Western Australia. The ISO code is listed as extinct. However, the Menang people remain in and around Albany, and the broader Noongar language continues to be spoken and taught. The Noongar article cited here gives fewer than 240 native speakers (2014 data) for Noongar as a whole. This report rests on a single search round. It found no dedicated Minang grammar, dictionary or school programme.

\section{Primary Educational Resources}
None specific to Minang were found.

\section{Linguistic and Reference Materials}
\begin{itemize}
\item Wikipedia, Noongar language (lists xrg under the Noongar group): \url{https://en.wikipedia.org/wiki/Noongar_language}
\item Wikipedia, Mineng (people and country): \url{https://en.wikipedia.org/wiki/Mineng}
\end{itemize}

\section{Religious and Community Resources}
\begin{itemize}
\item Historical Albany, ``King of Kinjarling'' (local history of Menang people): \url{https://historicalbany.com.au/king-of-kinjarling}
\item Historical Albany, Tommy King: \url{https://www.historicalbany.com.au/tommy-king}
\end{itemize}

\section{Digital and Audio Resources}
No Minang audio or digital resources were found. The City of Albany's dual-naming work with Menang Noongar elders on traditional place names is reported in an award listing: \url{https://www.infrastructure.gov.au/territories-regions-cities/local-government/awards/2021/promoting-indigenous-recognition-award-honourable-mention}

\section{Curriculum Development Context}
A Minang heritage curriculum would probably draw on wider Noongar teaching materials, which this search did not explore. It should be developed with Menang Noongar elders in Albany. Place names such as Kinjarling (Albany) offer a starting point for local content.

\end{document}